\section{Problem Setting}
\label{sec:problem}
% Rewritten 2026-10-29 after Adrian's comments (define "correct" first; the
% judging pass is a definition, not a method; first shift paragraph was
% related work). Previous version kept as comments at the end.
We study execution-free correctness prediction for code that an LLM generates
and judges itself, from signals of the model's own generation process: no tests are
executed and no additional reasoning tokens are generated. Given a programming
problem $x_i$, a model $\mathcal{M}$ samples candidate solutions
$c_{ij} \sim \mathcal{M}(x_i)$, $j = 1, \dots, n_i$. Following common practice in code generation benchmarks, a candidate is
considered \emph{correct} if it passes all tests. We denote the program correctness label as $y^{\mathrm{prog}}_{ij} \in \{0,1\}$.
Tests are executed only to obtain these labels for training and evaluation,
never to make a prediction. Where a statement holds for any problem and
candidate, we write $(x, c)$ and drop the indices.

\subsection{The Model's Verdict} % \orange{Heißt später "confidence".}}
\label{sec:problem-verdict}
To obtain the model's own judgment, we prompt $\mathcal{M}$ with the problem description $x_i$ and the candidate solution $c_{ij}$
and ask whether the implementation is correct, explicitely demanding a single token answer, \texttt{Yes} or \texttt{No}. We define the model's \emph{verdict} $v(x, c)$ as the
answer $t \in \{\texttt{Yes}, \texttt{No}\}$ with the higher probability
$p_{\mathcal{M}}(t {\mid} x, c)$. We refer to this probability as the {\it verdict confidence} and use it as a baseline in our experiments. The same forward pass that defines the verdict -- the \emph{judging
pass} --
 also generates the signals for the other prediction methods (Section~\ref{sec:methods}). We extract those at the verdict token's position: 
%Because the prompt ends in the same fixed
%text for every input, the verdict is produced at one position, the last
%prompt token, and the methods differ only in what they read there: the
probabilities of the two answer tokens {\tt Yes/No}, the hidden states, or a subspace of
the hidden states identified through interventions. 
%EigenScore and LAT use their own inputs.
%,
%see Section~\ref{sec:methods}.

\subsection{Prediction Targets}
We evaluate predictors for two different targets: % that capture different notions of correctness.

\paragraph{Program correctness ($y^{\mathrm{prog}}$)}
The first target is to predict whether the candidate implementation actually satisfies the problem specification, as operationalized by the test suite. This target corresponds to prior work on code correctness
prediction~\cite{bui2025openia, ribeiro2026repe, valentin2026incoherence,
sharma2025, he2026usingsemanticdistanceestimate}, and 
prediction scores are used to decide which candidates need testing.
Note that predictors may differ from the model's verdict $v$ as defined in Section \ref{sec:problem-verdict}, since they can be based on other internal signals. 

\paragraph{Verdict correctness ($y^{\mathrm{verd}}$)}
The second target is whether the model's verdict (as defined in Section \ref{sec:problem-verdict}) is correct, that is, whether
the model accepts a correct candidate or rejects an incorrect one:
\begin{equation*}
y^{\mathrm{verd}} = \mathbf{1}\bigl[v(x, c) = \texttt{Yes} \Leftrightarrow
y^{\mathrm{prog}} = 1\bigr].
\end{equation*}
A predictor for this target estimates when the model's judgment can be
trusted and when to invoke more expensive verification, such as test
execution or a stronger model. It is related to uncertainty estimation for generated answers, which predicts whether a model's own answer is
correct~\cite{Kadavath2022LanguageM, Huang2025Internal}.

Note that these two targets serve different purposes: program correctness concerns the
candidate itself, whereas verdict correctness concerns the reliability of the model's
judgment. When a model correctly rejects an incorrect candidate, for example,
its verdict is correct while the program is not. We evaluate both targets with the
same judging pass to investigate which internal signals support each prediction target.

\subsection{Distribution Shift}
\label{sec:problem-shift}
%A central challenge for execution-free correctness prediction is generalization beyond the training distribution, since a predictor may exploit regularities of its training data that do not hold
%elsewhere~\cite{OrgadTGRSKB25, Srey2026Factorised}. 
We study two
generalization scenarios, in which the evaluation data differ from the
training data in one controlled way.

In \emph{benchmark shift}, we train a predictor on one benchmark,
LiveCodeBench~\cite{Jain2025Livecodebench} or
BigCodeBench~\cite{Zhuo2025BCB}, and evaluate it on the other. This simulates
a predictor fitted on one kind of task and deployed on another:
LiveCodeBench consists of algorithmic programming challenges, BigCodeBench of
practical tasks that call libraries and APIs.

In \emph{language shift}, we train a predictor on candidates in one
programming language and evaluate it on candidates in another, for the same
problems in Python, C++ and Java. The problems are split once for all three
languages, so no problem is seen in training in one language and in
evaluation in another. 
%The shift thus changes the language of the code while
%the tasks stay comparable.

% ---------------------------------------------------------------------------
% Previous version (until 2026-10-28), condensed. The judging-prompt listing
% moved to appendix.tex (label app-listing:prompt).
%
% We study execution-free code correctness prediction as a means of steering
% agentic software development without repeatedly invoking tests or other
% verification tools. In particular, we investigate whether these signals
% generalize across benchmarks and programming languages, rather than relying
% on correlations specific to the training distribution. ... We construct a
% dataset D = {(x_i, c_ij, y_ij^prog)} over N problems, using test execution
% only to establish labels, not to make predictions. ... Such estimates could
% help prioritize candidates for verification or guide decisions within an
% agentic loop.
%
% Shared Judging Pass: We evaluate all methods on a common judging prompt ...
% The methods differ only in what they read there: the probabilities of the
% two answer tokens, the hidden states, or a subspace of the hidden states
% identified through interventions (Section III).
%
% Program correctness: ... A predictor for this target estimates whether the
% candidate is correct, regardless of the model's own assessment. Such a score
% can support candidate selection and help determine which generated programs
% require further verification.
% Verdict correctness: ... This target is related to uncertainty estimation for
% model-generated answers ... [Kadavath2022LanguageM, Huang2025Internal, TODO].
%
% Distribution Shift, first paragraph: A central challenge for execution-free
% correctness prediction is generalization beyond the training distribution.
% Probes on hidden states often fail to generalize across datasets
% [OrgadTGRSKB25, Levinstein2025LieDetector] and under production shifts
% [Kramar2026Probes]; for code, probe-based correctness assessors lose
% accuracy when applied to a model other than the one they were trained on
% [Vu2026Autoprobe]. A predictor may exploit regularities specific to the
% training problems, benchmark, or programming language rather than signals
% that reliably indicate correctness. Performance on held-out problems from
% the same benchmark therefore does not establish that a predictor will remain
% effective under distribution shift [Srey2026Factorised].
% In benchmark shift, the source and target domains consist of different
% benchmarks. ... In language shift, ... using a shared partition of the
% underlying problems to prevent overlap between training and evaluation
% problems. This setting tests whether correctness signals transfer across
% programming languages, where code representations and language-specific
% implementation details differ while the underlying problems correspond.
